\section*{Overview}

Persistent segment tree is the standard answer to ``keep all past versions.'' The useful mental model is not ``copy the
tree.'' It is ``copy only the nodes on the updated path and share everything else.''

\section*{Problem-Driven Motivation}

Many contest problems ask for queries on:

\begin{itemize}
  \item historical array versions,
  \item every prefix version of an array,
  \item old and new states interleaved in arbitrary order.
\end{itemize}

If we literally copy the whole segment tree per update, memory becomes \(O(nq)\). But a point update only affects
\(O(\log n)\) nodes. That mismatch is exactly where persistence pays off.

\section*{Recognition Pattern}

Persistent segment tree is a strong fit when:

\begin{itemize}
  \item updates are point updates or otherwise touch only one root-to-leaf path,
  \item queries refer to old versions,
  \item prefix versions or time snapshots are natural in the statement,
  \item an immutable version history is more useful than undo operations.
\end{itemize}

Typical contest signals are:

\begin{itemize}
  \item ``query version \(t\),''
  \item ``k-th smallest on subarray \([l,r]\)'' after building prefix versions,
  \item ``every update creates a new snapshot.''
\end{itemize}

\section*{Derivation}

In an ordinary segment tree, one point update changes exactly one root-to-leaf path. Every node outside that path still
has the correct aggregate.

So for a new version:

\begin{itemize}
  \item clone the nodes on the updated path,
  \item reuse all untouched children,
  \item return the new root pointer or index.
\end{itemize}

\includegraphics[width=\textwidth]{figures/version-branching.svg}

This is called path copying. It is the whole reason persistence is cheap here.

\section*{Worked Problem}

\paragraph{Problem.}
An array supports point assignments. After every update, keep the new version. Later queries ask for the sum on
\([l,r]\) in any version.

\paragraph{Why naive copying fails.}
Copying the full array or full segment tree after every update is too expensive in both time and memory.

\paragraph{How persistence fixes it.}
If version \(t\) updates position \(p\):

\begin{itemize}
  \item only the nodes whose segments contain \(p\) are cloned,
  \item all other subtrees are shared with version \(t-1\).
\end{itemize}

\paragraph{Why queries remain correct.}
Each root identifies one immutable tree whose aggregates match exactly the updates that were applied up to that version.

\paragraph{More advanced pattern.}
For k-th smallest on subarray \([l,r]\), build one persistent frequency tree per prefix. The difference between roots
\(\texttt{root[r]}\) and \(\texttt{root[l-1]}\) behaves like the multiset of the subarray.

\section*{Implementation Reasoning}

The implementation uses an index-based node pool rather than raw pointers. That makes cloning explicit:

\begin{itemize}
  \item \texttt{clone(node)} copies one old node into the pool,
  \item the new version root is the index returned by the top-level update,
  \item a vector of roots stores the full version history.
\end{itemize}

The invariant is simple:

\begin{quote}
Once created, a node is never mutated again by a later version unless that version first cloned it.
\end{quote}

That immutability is exactly what makes sharing safe.

\section*{Correctness Intuition}

Every version root points to a tree where each segment aggregate was recomputed from the correct children for that
version. Shared nodes correspond to untouched segments, so their aggregates remain valid. Cloned nodes cover exactly the
segments whose values changed.

\section*{Complexity Analysis}

\begin{itemize}
  \item one point update: \(O(\log n)\) time and \(O(\log n)\) new nodes,
  \item one query on any version: \(O(\log n)\),
  \item total memory after \(q\) updates: \(O(n + q \log n)\).
\end{itemize}

\section*{Common Pitfalls}

\begin{itemize}
  \item Accidentally mutating a shared node instead of cloning it.
  \item Forgetting to store every returned root.
  \item Underestimating the node-pool size when \(q\) is large.
  \item Using persistence when rollback or plain offline prefix processing would be simpler.
\end{itemize}

\section*{Variants and Failure Modes}

\begin{itemize}
  \item Persistent frequency trees are the standard order-statistics application.
  \item Persistent lazy propagation exists, but the implementation becomes much more fragile.
  \item If the task needs undo rather than arbitrary access to old versions, rollback structures are often simpler.
\end{itemize}

\section*{Practice Problems}

\begin{itemize}
  \item Historical point updates with range queries.
  \item K-th smallest on subarrays using prefix versions.
  \item Any offline query problem where every prefix snapshot must remain available.
\end{itemize}

\section*{References}

\begin{itemize}
  \item \href{https://cp-algorithms.com/data_structures/segment_tree.html}{cp-algorithms: Segment Tree, persistence section}.
  \item \href{https://github.com/kth-competitive-programming/kactl}{KACTL}.
  \item \href{https://codeforces.com/catalog}{Codeforces Catalog}.
\end{itemize}
